\documentclass[aspectratio=169,11pt]{beamer}
\usepackage{../../../shared/theme/esmad_beamer_theme}
\usepackage{amsmath,amssymb}
\usepackage{booktabs}

\title[HTE and Causal Forests]{Heterogeneous Treatment Effects and Causal Forests}
\subtitle{From Average Effects to Conditional Effects}
\date{\today}

\newcommand{\E}{\mathbb{E}}

\begin{document}

\begin{frame}
\titlepage
\end{frame}

\begin{frame}{Learning Goals}
\begin{itemize}
\item distinguish ATE, subgroup effects, and conditional average treatment effects;
\item explain why predicting outcomes is different from predicting treatment effects;
\item understand causal trees and causal forests;
\item explain honesty, sample splitting, and orthogonalization;
\item diagnose overlap and unstable subgroup discoveries;
\item connect heterogeneity estimates to treatment targeting without overstating precision.
\end{itemize}
\end{frame}

\begin{frame}{Outline}
\tableofcontents
\end{frame}

\section{From Average to Conditional Effects}

\begin{frame}{Average Treatment Effect}
The ATE is
\[
ATE=\E[Y(1)-Y(0)].
\]

It summarizes the population average, but it can hide important variation across units.
\end{frame}

\begin{frame}{Conditional Average Treatment Effect}
For covariates \(X=x\),
\[
\tau(x)=\E[Y(1)-Y(0)\mid X=x].
\]

This is the conditional average treatment effect, or CATE.

\begin{block}{Question}
Which observed characteristics modify the causal response to treatment?
\end{block}
\end{frame}

\begin{frame}{Effect Modification vs Prognosis}
A variable can strongly predict outcomes without modifying treatment effects.

For example, baseline risk may affect \(Y\) but not the difference
\[
Y(1)-Y(0).
\]

\begin{alertblock}{Key distinction}
Prognostic importance is not causal-effect heterogeneity.
\end{alertblock}
\end{frame}

\section{Why Heterogeneity Is Hard}

\begin{frame}{The Individual Effect Is Never Observed}
For each unit,
\[
\tau_i=Y_i(1)-Y_i(0),
\]
but only one potential outcome is observed.

Therefore there is no direct supervised-learning target for individual treatment effects.
\end{frame}

\begin{frame}{Identification Still Comes First}
Under conditional exchangeability,
\[
(Y(1),Y(0))\perp D\mid X,
\]
plus overlap and consistency, CATEs can be identified.

Flexible ML does not repair unmeasured confounding or lack of support.
\end{frame}

\section{Causal Trees}

\begin{frame}{Recursive Partitioning for Treatment Effects}
A causal tree splits the covariate space to create groups with different estimated treatment effects.

Unlike a prediction tree, the splitting criterion targets heterogeneity in causal effects rather than outcome fit.
\end{frame}

\begin{frame}{Honest Estimation}
Honest trees separate:
\begin{itemize}
\item data used to choose splits;
\item data used to estimate treatment effects within leaves.
\end{itemize}

\begin{block}{Purpose}
Reduce adaptive bias caused by using the same noise both to discover and estimate subgroups.
\end{block}
\end{frame}

\section{Causal Forests}

\begin{frame}{From Trees to Forests}
Causal forests average many causal trees.

The goals are:
\begin{itemize}
\item flexible CATE estimation;
\item reduced instability relative to a single tree;
\item local weighting of observations with similar covariates.
\end{itemize}
\end{frame}

\begin{frame}{Local Weighting View}
A forest estimate can be viewed schematically as
\[
\widehat\tau(x)
=
\sum_i \alpha_i(x)\widehat\psi_i,
\]
where the forest determines local weights \(\alpha_i(x)\) and \(\widehat\psi_i\) is an orthogonalized treatment-effect score.
\end{frame}

\section{Orthogonalization}

\begin{frame}{Nuisance Functions}
Define
\[
m(x)=\E[Y\mid X=x]
\]
and
\[
e(x)=P(D=1\mid X=x).
\]

Orthogonalization residualizes outcome and treatment using these nuisance models before estimating heterogeneity.
\end{frame}

\begin{frame}{Robinson-Style Residualization}
A useful representation is
\[
Y-m(X)
=
\tau(X)(D-e(X))
+
\varepsilon.
\]

This reduces sensitivity to errors in nuisance estimation under regularity conditions.
\end{frame}

\begin{frame}{Cross-Fitting}
Estimate nuisance models on one fold and evaluate residualized scores on another.

Cross-fitting helps avoid reusing the same observations for flexible nuisance fitting and effect estimation.
\end{frame}

\section{Overlap and Support}

\begin{frame}{Heterogeneity Requires Local Overlap}
For CATE estimation, overlap must hold within relevant covariate regions:
\[
0<e(x)<1.
\]

If a subgroup is almost always treated or untreated, its treatment effect is largely extrapolation.
\end{frame}

\begin{frame}{Extreme-CATE Warnings}
Very large estimated CATEs can arise from:
\begin{itemize}
\item weak overlap;
\item small effective samples;
\item noisy leaves;
\item adaptive subgroup discovery.
\end{itemize}

Always inspect support before interpreting extremes.
\end{frame}

\section{Validation}

\begin{frame}{How Do We Validate Heterogeneity?}
We cannot compare estimated individual effects to observed truth.

Instead evaluate:
\begin{itemize}
\item calibration of predicted treatment effects;
\item group-average effects across predicted-effect quantiles;
\item out-of-sample policy value;
\item stability across samples and specifications.
\end{itemize}
\end{frame}

\begin{frame}{Sorted Group Effects}
Rank units by predicted CATE and form groups.

Estimate treatment effects within those groups using held-out or orthogonalized procedures.

\begin{block}{Goal}
Check whether units predicted to benefit more actually show larger group-average effects.
\end{block}
\end{frame}

\begin{frame}{Subgroup Multiplicity}
Searching many covariates creates opportunities for spurious subgroup findings.

\begin{itemize}
\item pre-specify important subgroups where possible;
\item use honest estimation;
\item validate on held-out data;
\item avoid storytelling around one unstable split.
\end{itemize}
\end{frame}

\section{Policy Targeting}

\begin{frame}{From CATE to Treatment Rules}
A simple treatment rule may be
\[
\pi(x)=\mathbb{1}\{\widehat\tau(x)>c\},
\]
where \(c\) reflects treatment cost or a decision threshold.
\end{frame}

\begin{frame}{Policy Value}
The relevant target is often not CATE prediction error but expected outcome under the learned policy:
\[
V(\pi)=\E[Y(\pi(X))].
\]

\begin{alertblock}{Decision focus}
A slightly imperfect CATE model can still yield a good policy, while a visually complex heterogeneity model may yield little decision value.
\end{alertblock}
\end{frame}

\section{Applied Reasoning}

\begin{frame}{Example: Regional Pricing Intervention}
Suppose the average pricing-policy effect is negative but heterogeneous.

Potential modifiers:
\begin{itemize}
\item baseline demand;
\item regional income;
\item competitor intensity;
\item customer tenure;
\item seasonality exposure.
\end{itemize}

The question becomes which regions are most responsive, not just whether the average effect is nonzero.
\end{frame}

\begin{frame}{A Credible HTE Workflow}
\begin{enumerate}
\item Define the CATE or subgroup estimand.
\item Establish causal identification first.
\item Fit nuisance models with cross-fitting.
\item Estimate heterogeneity honestly.
\item Inspect overlap and effective sample sizes.
\item Validate heterogeneity out of sample.
\item Compare group-average effects across predicted CATE strata.
\item Evaluate policy value if targeting is the goal.
\item Report instability and uncertainty prominently.
\end{enumerate}
\end{frame}

\begin{frame}{Common Failure Modes}
\begin{itemize}
\item confusing outcome prediction with treatment-effect prediction;
\item interpreting noisy individual CATEs as known personal effects;
\item searching many subgroups without validation;
\item ignoring weak overlap;
\item using the same data for discovery and confirmation;
\item treating variable importance as causal-effect importance;
\item optimizing heterogeneity without checking policy value.
\end{itemize}
\end{frame}

\begin{frame}{Synthesis: Heterogeneity Without Storytelling}
\begin{itemize}
\item CATEs describe how average causal effects vary with observed covariates.
\item Individual causal effects remain fundamentally unobserved.
\item Causal trees target treatment-effect differences rather than outcome prediction.
\item Honest forests and cross-fitting reduce adaptive bias.
\item Orthogonalization reduces first-order sensitivity to nuisance estimation.
\item Local overlap is essential for credible subgroup effects.
\item Heterogeneity should be validated out of sample and connected to explicit decisions.
\end{itemize}
\end{frame}

\begin{frame}{Selected References}
\small
\begin{itemize}
\item Athey, S., and Imbens, G. (2016). Recursive Partitioning for Heterogeneous Causal Effects. \textit{PNAS}, 113(27), 7353--7360.
\item Wager, S., and Athey, S. (2018). Estimation and Inference of Heterogeneous Treatment Effects Using Random Forests. \textit{JASA}, 113(523), 1228--1242.
\item Athey, S., Tibshirani, J., and Wager, S. (2019). Generalized Random Forests. \textit{Annals of Statistics}, 47(2), 1148--1178.
\item Chernozhukov, V. et al. (2018). Double/Debiased Machine Learning for Treatment and Structural Parameters. \textit{The Econometrics Journal}, 21(1), C1--C68.
\item Nie, X., and Wager, S. (2021). Quasi-Oracle Estimation of Heterogeneous Treatment Effects. \textit{Biometrika}, 108(2), 299--319.
\end{itemize}
\end{frame}

\contactslide

\end{document}
